\newpage
\chapter{Courses}
\begin{articles}
  \item This chapter establishes the necessary provisions concerning the conduct of courses, in accordance with the principles of the Charter.
  \item The Academy has a Rhythm Games Foundations Faculty and a Rhythm Games Practice Faculty.
  \begin{paragraphs}
    \item The course categories within each Faculty are as follows.
  \end{paragraphs}
  \begin{subitems}
    \item Rhythm Games Foundations Faculty: Creation; Humanities and Sciences.
    \item Rhythm Games Practice Faculty: Arcade; Standalone; Mobile.
  \end{subitems}
  \item The Rhythm Games Foundations Faculty covers knowledge, culture, and creation relating to rhythm games.
  \begin{subitems}
    \item Creation: fields involving art, design, development, and other forms of creation.
    \item Humanities and Sciences: fields within the humanities, social sciences, and natural sciences.
  \end{subitems}
  \item The Rhythm Games Practice Faculty covers skills and knowledge relating to playing rhythm games.
  \begin{subitems}
    \item Arcade: fields relating to rhythm games played on cabinets installed in arcades and other premises.
    \item Standalone: fields relating to rhythm games running on home consoles, computers, and other equipment.
    \item Mobile: fields relating to rhythm games running on mobile devices.
  \end{subitems}
  \item The Faculties and course categories specified in the preceding two Articles are guides for finding courses and do not limit the scope of Students' activities.
  \begin{paragraphs}
    \item For a course spanning several categories, the principal category is selected.
    \item For a course that is difficult to place in any category, the closest category is selected. In such cases, Deans provide advice on choosing a category.
  \end{paragraphs}
  \item A Lecturer may open a course by submitting the prescribed application and receiving approval.
  \begin{paragraphs}
    \item The approval referred to in the preceding paragraph requires the unanimous agreement of the Deans.
    \item The specific application procedure referred to in paragraph 1 is established separately.
  \end{paragraphs}
  \item Students may take courses in accordance with separately established provisions.
  \begin{paragraphs}
    \item A Lecturer may set conditions for taking their course. In that case, the Lecturer must make those conditions known in advance to Students wishing to take it.
    \item Where course registration is required to take a course, the procedures concerning the method of registration, when it takes effect, changes, and cancellation are established separately.
    \item Access to course registration and participation records is limited to the Lecturer responsible for the course and Deans who need access for administrative work. Use of those records is limited to what is necessary to identify participants, change or cancel registrations, communicate as required to run the course, and enable the Lecturer to record participation and assessments.
    \item Course registration and assessment records kept by a Lecturer do not constitute the Academy's formal recognition of grades or credits. The system for grades and credits is established separately.
  \end{paragraphs}
  \item Copyright and other intellectual property rights in course materials belong to the Lecturer who created them.
  \begin{paragraphs}
    \item Notwithstanding the preceding paragraph, on the basis of the consent referred to in Chapter 6, Article 2, paragraph 1, item (1), the Academy may use and distribute course materials free of charge to the extent necessary for activities within the Academy.
  \end{paragraphs}
  \item If a Lecturer has difficulty conducting the planned number of course sessions during the regular activity period, the remaining sessions may be conducted during the make-up period.
  \item Deans may ask a Lecturer to discontinue or change a course only when they determine that facts falling within one of the following items exist.
  \begin{subitems}
    \item The communication of clearly false information or other facts that seriously undermine the integrity of the course.
    \item Harassment or other seriously improper conduct toward participants or other Students.
    \item In addition to the preceding two items, facts that are considered to pose a risk of seriously harming Students' rights or interests.
  \end{subitems}
  \begin{paragraphs}
    \item A request under the preceding paragraph must be limited to the minimum necessary so as not to unduly restrict the academic freedom provided for in the Charter.
  \end{paragraphs}
  \item The Deans separately establish provisions needed to implement this chapter.
\end{articles}
